
\begin{center}
    {\Large\itshape Alban (au temps des Croisades : XII\textsuperscript{ème} siècle)}
\end{center}

\vspace{1em}

Pour protéger, défendre les « lieux saints » ou les pèlerins s'y
rendant, se fondent des ordres militaires -- Hospitaliers --
Templiers. Ces moines-soldats s'établissent chez nous autant qu'on
puisse le conclure par des documents écrits et des noms de lieux.

\vspace{1em}

Ainsi, le 11 février 1161, Raymond Trencavel donne aux chanoines
de Sainte Cécile l'église et l'hôpital de Carmé (ce qui signifie
les biens des hospitaliers de Montfranc (St Léonard del Carmé).
En 1217, à la suite de difficultés entre le prieur et les frères
de la MAIO de Rayssac avec les hospitaliers du Carmé (Montfranc)
et les chanoines de Ste Cécile, ce fut le prévôt de Vieux et
Bernat de St Sernin qui tranchèrent le différend. En 1447, le
commandeur de Rayssac vend des terres sises à Négremontel
(appartenant aux hospitaliers de Carmésal) à R. Limorin d'Alban.
En 1598, le compoix d'Alban y note la présence d'hospitaliers et
dresse la nomenclature de leurs biens. Un acte du 10 décembre
1682 signale un hospitalier en la personne de Pierre Clergué,
sieur de la Bouyssounade.
